\chapter{System Architecture and Design}\label{ch:ch04}

This chapter describes how APEX is put together. It starts from the outermost
view --- the running processes and the traffic that flows between them --- and
then works inward: the lifecycle a single HTTP request follows through the
backend, the reasoning behind each technology choice, the shape of the public
API, the physical layout of the source tree, and the design patterns that hold
the whole thing together. The guiding principle throughout is deliberate
simplicity: APEX is a graduation project intended to be self-hosted by a single
instructor or a small department, so every architectural decision is weighed
against the cost of operating it, not against the demands of hyperscale.

\section{Architectural Overview}\label{sec:arch-overview}

APEX is a conventional client--server web application built around a
single-page application (SPA) front end and a stateless HTTP API back end that
persists everything of value to a relational database. What distinguishes it
from a textbook three-tier system is the deliberate refusal to multiply moving
parts: there is exactly one server process to run, one database to back up, and
one external dependency that actually executes untrusted code. Everything else
is optional.

\Cref{fig:architecture} shows the full runtime and deployment topology. Six
components appear, and it is worth being precise about which of them are
first-class services and which are external dependencies, because the
distinction drives both the security posture and the operational story.

\begin{figure}[htbp]\centering
\resizebox{\textwidth}{!}{%
\begin{tikzpicture}[font=\sffamily\small]
  % ---- Self-hosted (docker-compose) tier --------------------------------
  \node[apexbox] (spa) {React + Vite SPA\\[1pt]\scriptsize served by nginx};
  \node[apexbox, right=34mm of spa] (api)
    {Gin REST API\\[1pt]\scriptsize \texttt{/api/*}\ \ \texttt{/healthz}};

  % in-process reminder goroutine attached to the API
  \node[draw=apexblue, dashed, fill=apexblue!6, rounded corners=2pt,
        align=center, font=\sffamily\scriptsize, above=9mm of api] (rem)
    {reminder goroutine\\(1-minute tick)};

  % PostgreSQL store below the API
  \node[apexstore, below=16mm of api] (db)
    {PostgreSQL 16\\[1pt]\scriptsize accessed via GORM};

  % side migration CLI feeding the database
  \node[apexbox, left=24mm of db] (mig)
    {\texttt{cmd/migrate}\\[1pt]\scriptsize golang-migrate};

  % ---- External services (NOT compose services) -------------------------
  \node[apexext, right=42mm of api, yshift=12mm] (judge)
    {Judge0 CE (external)\\[1pt]\scriptsize inner sandbox: \texttt{isolate}};
  \node[apexext, right=42mm of api, yshift=-14mm] (smtp)
    {SMTP (external, optional)\\[1pt]\scriptsize password reset + reminders};

  % ---- Edges ------------------------------------------------------------
  \draw[apexarrowblue] (spa) -- node[apexlabel, above]{HTTPS / JSON} (api);
  \draw[apexarrow] (api) -- node[apexlabel, right]{SQL} (db);
  \draw[apexarrow] (mig) -- node[apexlabel, above]{SQL} (db);
  \draw[apexarrow] (api.east) -- node[apexlabel, above, sloped]{HTTP (\texttt{JUDGE0\_URL})} (judge.west);
  \draw[apexarrow] (api.east) -- node[apexlabel, below, sloped]{SMTP} (smtp.west);
  \draw[apexarrowblue, dashed] (rem) -- (api);

  % ---- Deployment boundaries (background layer) -------------------------
  \begin{scope}[on background layer]
    \node[draw=apexgrey, dashed, rounded corners=3pt, inner sep=6mm,
          fit=(spa)(api)(rem)(db)(mig),
          label={[font=\sffamily\scriptsize\itshape, anchor=north west]north west:%
            Self-hosted \texttt{docker-compose} tier}] (compose) {};
    \node[draw=apexamber, dashed, rounded corners=3pt, inner sep=4mm,
          fit=(judge)(smtp),
          label={[font=\sffamily\scriptsize\itshape, anchor=north east]north east:%
            External (not compose services)}] {};
  \end{scope}
\end{tikzpicture}%
}
\caption{APEX runtime and deployment architecture.}\label{fig:architecture}
\end{figure}


\begin{itemize}
  \item \textbf{React + Vite SPA.} The browser-side application, compiled to a
        static bundle and served by \texttt{nginx} in production. It holds no
        durable state of its own beyond a JWT in \texttt{localStorage}; every
        meaningful action is a JSON call to the API.
  \item \textbf{Gin REST API.} The single Go process (binary \texttt{apex})
        that owns all business logic. It exposes routes under \texttt{/api/*}
        plus an unauthenticated \texttt{/healthz} liveness endpoint. This
        process is the heart of the system and the subject of most of this
        chapter.
  \item \textbf{PostgreSQL 16.} The single source of durable truth, reached
        exclusively through the GORM object--relational mapper. No component
        other than the API and the migration tool ever touches it directly.
  \item \textbf{Judge0 CE (external).} The sandboxed code-execution service
        that compiles and runs student submissions inside the \texttt{isolate}
        jail. Crucially, Judge0 is \emph{not} a \texttt{docker-compose} service;
        it is reached over HTTP at the URL named by \texttt{JUDGE0\_URL}
        (defaulting to the public \texttt{https://ce.judge0.com}). Treating the
        single most security-sensitive component --- the one that runs
        arbitrary code --- as an external service reached through one URL is an
        intentional isolation boundary.
  \item \textbf{SMTP relay (external, optional).} Used for password-reset links
        and exam reminders. It degrades gracefully: when \texttt{SMTP\_HOST} is
        empty, every message that would have been sent is logged to standard
        output instead, so the system remains fully functional without a mail
        provider configured.
  \item \textbf{\texttt{cmd/migrate}.} A separate one-shot binary, built from
        the same Go module, that owns schema evolution in production. It applies
        versioned SQL migrations and exits.
\end{itemize}

Two further behaviours live \emph{inside} the API process rather than as
separate services, and \cref{fig:architecture} draws them attached to it. The
first is the \textbf{reminder goroutine}, a background loop that wakes on a
one-minute ticker to dispatch exam reminders. The second, not drawn but equally
important, is the fleet of \textbf{fire-and-forget grading goroutines} spawned
whenever an exam is submitted. Both are examples of a recurring theme: work
that a larger system would hand to a dedicated worker or cron daemon is folded
into the single binary, because for a classroom-scale tool the operational
saving of ``one process to run and one process to reason about'' outweighs the
resilience a separate worker fleet would buy.

The deployment boundary in \cref{fig:architecture} is significant. The
self-hosted tier --- SPA, API, database, and the migration tool --- is what a
\texttt{docker-compose up} brings up on a single host. The external tier ---
Judge0 and SMTP --- is reached over the network and is never bundled. This
separation means that an operator can run the entire teaching platform on one
modest machine while pointing at a shared or public Judge0 instance, and can
add email only when they are ready to configure a relay.

\section{The Four-Layer Request Lifecycle}\label{sec:lifecycle}

Every request that reaches the backend passes through the same four conceptual
layers, in the same order. This uniformity is what keeps a six-thousand-line Go
codebase navigable: once the layering is understood, any endpoint can be read
by asking which layer each line belongs to. \Cref{fig:layers} traces the
pipeline from an incoming HTTP request to a JSON response, including the two
short-circuit paths where a request is rejected before it reaches the handler.

\begin{figure}[htbp]
\centering
\resizebox{\textwidth}{!}{%
\begin{tikzpicture}[
    node distance=10mm,
    layer/.style={apexbox, text width=32mm, minimum height=30mm,
      align=center, font=\sffamily\scriptsize, inner sep=5pt},
    tag/.style={apexlabel, draw=apexgrey, rounded corners=2pt, fill=apexlight,
      minimum height=10mm, text width=18mm, align=center, font=\sffamily\scriptsize},
  ]
  \def\hd#1{{\sffamily\small\bfseries\color{apexblue}#1}}
  \def\it#1{{\sffamily\scriptsize\color{apexgrey!25!black}#1}}

  % Incoming request
  \node[tag] (req) {HTTP\\request};

  % Layer 1 -- Router
  \node[layer, right=of req] (router)
    {\hd{1. Router}\\[4pt]
     \it{public vs.\\ protected groups}\\[3pt]
     \it{\texttt{/api} dispatch}};

  % Layer 2 -- Middleware
  \node[layer, right=of router] (mw)
    {\hd{2. Middleware}\\[4pt]
     \it{CORS allowlist}\\[3pt]
     \it{JWT authenticate}\\[3pt]
     \it{RequireRole}};

  % Layer 3 -- Handler
  \node[layer, right=of mw] (handler)
    {\hd{3. Handler}\\[4pt]
     \it{bind + validate}\\[3pt]
     \it{ownership check}\\[3pt]
     \it{business logic}};

  % Layer 4 -- Data (two backing stores)
  \node[apexstore, right=of handler, text width=32mm, minimum height=13mm,
        align=center, font=\sffamily\scriptsize] (db)
    {\hd{4. Data}\\[3pt]
     \it{GORM / PostgreSQL}};
  \node[apexext, below=6mm of db, text width=32mm, minimum height=11mm,
        align=center, font=\sffamily\scriptsize] (j0)
    {\it{Judge0 client (\texttt{isolate})}};

  % Response
  \node[tag, below=6mm of router] (resp) {JSON\\response};

  % Forward pipeline arrows
  \draw[apexarrowblue] (req) -- (router);
  \draw[apexarrowblue] (router) -- (mw);
  \draw[apexarrowblue] (mw) -- (handler);
  \draw[apexarrowblue] (handler) -- (db);
  \draw[apexarrow] (handler.south) |- (j0.west);

  % Rejection / return path
  \draw[apexarrow] (mw.south) |- (resp)
    node[apexlabel, pos=0.25] {401 / 403};
  \draw[apexarrow] (db.south) |- ([yshift=-1mm]resp.east);

\end{tikzpicture}%
}
\caption{Four-layer request lifecycle through the Gin backend.}
\label{fig:layers}
\end{figure}


\subsection{Layer 1 --- Router and route groups}

The transport layer is assembled in \texttt{main.go}. After the process
initialises its configuration, database connection, and subsystems, it builds a
Gin engine and installs global middleware:

\begin{lstlisting}[language=Go,caption={Router construction and route-group split in \texttt{main.go}.},label={lst:router}]
r := gin.Default()
r.Use(middleware.CORS(cfg.AllowedOrigins))

r.GET("/healthz", func(c *gin.Context) {
    c.JSON(http.StatusOK, gin.H{"status": "ok"})
})

public := r.Group("/api")
protected := r.Group("/api")
protected.Use(middleware.Auth(cfg.JWTSecret))
\end{lstlisting}

The design pivots on the two route groups. Both share the \texttt{/api} prefix,
but \texttt{public} carries no authentication middleware while \texttt{protected}
applies \texttt{middleware.Auth} to \emph{every} route registered on it. This
single \texttt{.Use} call is the mechanism that separates ``anyone may call
this'' (login, signup, password reset) from ``you must present a valid token.''
The \texttt{/healthz} endpoint is registered directly on the engine, outside
both groups, so a load balancer or the container runtime can probe liveness
without a token; it returns \texttt{\{"status":"ok"\}} and deliberately touches
neither the database nor Judge0, making it a pure ``is the process routing?''
check.

\subsection{Layer 2 --- Middleware}

Middleware runs before any handler and is where cross-cutting concerns are
enforced. For a protected request the chain is, from outermost to innermost: the
Gin \texttt{Logger} and \texttt{Recovery} middleware supplied by
\texttt{gin.Default()} (the latter converts a handler panic into a 500 rather
than crashing the process); the CORS filter; the JWT authentication gate; and,
where a route or subgroup demands it, a role guard. Because rejection happens
here, a malformed or unauthorised request is turned away with a 401 or 403
\emph{before} it can reach business logic or the database --- the left-hand
short-circuit arrow in \cref{fig:layers}.

The CORS middleware is intentionally strict. It allows only the origins
assembled at startup (\texttt{APP\_URL} plus any comma-separated
\texttt{ALLOWED\_ORIGINS}), only the methods \texttt{GET}, \texttt{POST},
\texttt{PUT}, \texttt{DELETE}, and \texttt{OPTIONS}, and only the
\texttt{Content-Type} and \texttt{Authorization} headers. It sets
\texttt{AllowCredentials: false}, which is a considered decision: APEX carries
its token in the \texttt{Authorization} header rather than in a cookie, so there
are no ambient credentials for a cross-site request to abuse, and preflight
results are cached for twelve hours to cut down on \texttt{OPTIONS} chatter.

\subsection{Layer 3 --- Handler}

Once a request clears the middleware, the handler runs the actual work. Handlers
across the codebase follow the same skeleton: read the verified identity from
the Gin context (\texttt{c.GetUint("user\_id")}), read any path parameters, load
the target resource and check that the caller is entitled to it, bind and
validate the JSON body with \texttt{c.ShouldBindJSON}, perform the database work
through GORM, trigger any background work, and finally serialise a JSON response
with an appropriate status code. Handlers never trust an identifier taken from
the request body when deciding ownership; they use only the \texttt{user\_id}
that the authentication middleware placed on the context.

\subsection{Layer 4 --- Data}

The innermost layer is persistence. In the common case this means the shared
\texttt{database.DB} handle --- a \texttt{*gorm.DB} configured with a bounded
connection pool (25 open connections, 10 idle, a five-minute idle timeout and a
thirty-minute maximum lifetime) --- which GORM translates into SQL against
PostgreSQL. For coding submissions a second backing store appears at this layer:
the Judge0 client, which sends code over HTTP to the sandbox. \Cref{fig:layers}
draws this as a branch off the handler, because a request that grades code
touches both PostgreSQL and Judge0 before it can respond.

\subsection{An honest note on the layering}

It is worth being candid about what this layering is and is not. It is a
pragmatic, package-level separation of transport, handler, service, and
persistence concerns. It is \emph{not} a ports-and-adapters or ``clean''
architecture: there is no repository interface abstracting the database, and
handlers reach directly into the global \texttt{database.DB}. The service layer
is simply the set of packages (\texttt{judge0}, \texttt{notification},
\texttt{email}, \texttt{reminder}) that handlers call into. For a single-binary
classroom tool this trade favours readability and low ceremony over the
testability and swappability that dependency inversion would provide, and the
book states that plainly rather than dressing it up.

\section{Three-Gate Authorization}

Layering the request lifecycle is only half of the access-control story: on top
of it APEX applies three independent, fail-closed authorization gates ---
authentication (\texttt{middleware.Auth}), a role check
(\texttt{middleware.RequireRole}), and a resource-ownership check inside each
handler. Because these gates belong conceptually to the request pipeline they
are named here at an architectural level, but their full treatment --- how each
one fails closed and why a bug in one cannot open the others --- lives in
\cref{ch:ch08} and is depicted in \cref{fig:authz}.

\section{Technology Stack and Rationale}\label{sec:stack}

Every stack choice in APEX was made to keep the operational surface small and
the code readable by a five-person team. \Cref{tab:stack} summarises the
selections; the subsections that follow justify the load-bearing ones.

\begin{table}[htbp]
\centering
\caption{Principal technology choices and their roles.}\label{tab:stack}
\begin{tabularx}{\linewidth}{@{}L{3.2cm}L{4.3cm}X@{}}
\toprule
\textbf{Concern} & \textbf{Choice} & \textbf{Role in APEX} \\
\midrule
Backend language & Go 1.25.3 & Compiles to one static binary; simple concurrency \\
HTTP framework & \texttt{gin-gonic/gin} 1.12 & Routing, middleware, JSON binding \\
ORM & \texttt{gorm.io/gorm} 1.31 & Struct-to-table mapping, query building \\
Database & PostgreSQL 16 & Durable relational store \\
Sandbox & Judge0 CE & Isolated execution of untrusted code \\
Front end & React 18 + Vite 5 + TypeScript 5 & Typed single-page application \\
Migrations & \texttt{golang-migrate/migrate/v4} & Versioned schema evolution in production \\
\bottomrule
\end{tabularx}
\end{table}

\subsection{Why Go and Gin}

Go~\cite{go} was chosen for a property that shows up repeatedly in the
deployment story of \cref{ch:ch10}: it compiles the entire backend into a
single, self-contained, statically linked binary. There is no interpreter to
install, no runtime to match, and no dependency tree to reconstruct on the
server --- the artefact is copied and run. Go's goroutines also make the
concurrency that APEX needs (background grading, the reminder ticker)
inexpensive to express: a background task is literally the keyword \texttt{go}
in front of a function call.

Gin~\cite{gin} sits on top of Go's standard \texttt{net/http} and supplies
exactly the pieces a JSON API needs without imposing a heavyweight framework:
fast radix-tree routing, a composable middleware chain, route groups, and
convenience helpers for binding and emitting JSON. The route-group feature in
particular is what makes the clean public-versus-protected split of
\cref{lst:router} possible in three lines. Gin does not dictate a project
structure, which let the team impose the package layout of
\cref{sec:packages} rather than fight a framework's conventions.

\subsection{Why GORM}

GORM~\cite{gorm} is the most widely used object--relational mapper in the Go
ecosystem, and it earns its place here by removing the tedium of hand-writing
SQL for fifteen entities while still allowing raw SQL where it matters. Handlers
express intent in Go --- \texttt{database.DB.Where("id = ? AND teacher\_id = ?",
examID, teacherID).First(\&exam)} --- and GORM generates the parameterised
query, which as a side benefit closes the door on SQL injection for every
query written this way. GORM's \texttt{AutoMigrate} facility, which reflects
over the model structs to create tables, accelerates development enormously.
That same facility is also the origin of the project's most instructive bug ---
the \texttt{is\_draft} data-loss incident detailed in \cref{ch:ch05} --- which
is precisely why production disables \texttt{AutoMigrate} and hands schema
authority to the versioned migration binary instead.

\subsection{Why PostgreSQL}

PostgreSQL~\cite{postgres} is the durable store. The data is inherently
relational --- users own classes, classes contain members, exams link to
classes, problems belong to exams, submissions belong to attempts --- and a
relational database with real foreign keys and unique constraints models this
faithfully. Three composite unique indexes in particular
(\texttt{ClassMember(class\_id, user\_id)},
\texttt{ExamAttempt(user\_id, exam\_id)}, and
\texttt{ExamClass(exam\_id, class\_id)}) enforce important invariants --- a
student joins a class at most once, and holds at most one attempt row per exam
--- at the database level, so that even a race between concurrent requests
cannot violate them. PostgreSQL's JSONB columns additionally give APEX a place
to store semi-structured data (MCQ options, autosave payloads, announcement
attachments) without abandoning the relational model for the fields that
benefit from schema. Version 16 runs identically in a local Docker container
and on managed providers such as Neon or Railway, reached through a single
\texttt{DATABASE\_URL}.

\subsection{Why React, Vite, and TypeScript}

The front end is a React~\cite{react} single-page application written in
TypeScript~\cite{typescript} and built with Vite~\cite{vite}. React's
component model suits an application that is essentially many stateful views ---
an exam builder, a Monaco-based coding view, a grading queue --- reused across
routes. TypeScript adds static types across roughly seventeen thousand lines of
front-end code, catching whole classes of shape mismatches between the client
and the JSON contracts described in \cref{sec:api} before they ever reach the
browser. Vite provides a fast development server with hot module replacement and
a production build that emits content-hashed, aggressively cacheable static
assets. The result of the build is a folder of static files with no server-side
rendering, which is why it can be served by a stock \texttt{nginx} container.

\subsection{Why one binary that also ships \texttt{migrate} and \texttt{seed}}

A single Go module produces three binaries: \texttt{apex} (the long-running API
server and in-process scheduler), \texttt{migrate} (a one-shot schema tool), and
\texttt{seed} (an on-demand demo-data loader). Shipping all three from one
module, and baking all three into one Docker image, means there is exactly one
build to trust and one set of dependencies to audit. The production container
runs \texttt{apex} by default but the compose stack overrides its entrypoint to
run \texttt{migrate up} for the one-shot migration service, so the same image
serves two roles. This is the same economy of ``one artefact, one thing to
reason about'' that motivated the choice of a compiled single-binary language in
the first place, applied at the level of the deployment.

\section{API Design}\label{sec:api}

APEX exposes a RESTful HTTP API under the base path \texttt{/api}, following the
resource-oriented style described by Fielding~\cite{fielding2000rest}. Nouns
name resources (\texttt{/exams}, \texttt{/classes}, \texttt{/submissions}), HTTP
verbs express intent (\texttt{GET} to read, \texttt{POST} to create,
\texttt{PUT} to update, \texttt{DELETE} to remove), and every request and
response body is JSON. Authentication is stateless: a bearer token issued at
login is presented on each protected call, so the server keeps no session state
and any instance can serve any request.

\subsection{Modular route registration}

Rather than declaring every route in one place, \texttt{main.go} delegates route
registration to thirteen feature modules, each of which exposes a
\texttt{RegisterRoutes(public, protected)} function and decides for itself which
group and which role guard each of its endpoints needs:

\begin{lstlisting}[language=Go,caption={Modular route registration in \texttt{main.go}.},label={lst:register}]
auth.RegisterRoutes(public, protected)
class.RegisterRoutes(public, protected)
student.RegisterRoutes(public, protected)
exam.RegisterRoutes(public, protected)
problem.RegisterRoutes(public, protected)
testcase.RegisterRoutes(public, protected)
submission.RegisterRoutes(public, protected)
execute.RegisterRoutes(public, protected)
notification.RegisterRoutes(public, protected)
profile.RegisterRoutes(public, protected)
teacher.RegisterRoutes(public, protected)
announcement.RegisterRoutes(public, protected)
folder.RegisterRoutes(public, protected)
\end{lstlisting}

The only thing each module receives is the two router groups. There is no
dependency-injection container and no service registry; the ``wiring'' is
literally the passing of two \texttt{*gin.RouterGroup} values. This is a
conscious choice in favour of explicitness --- to learn what routes exist, one
reads thirteen small \texttt{routes.go} files rather than tracing a
configuration graph.

\subsection{Public, protected, and role subgroups}

Modules apply access control at three granularities, all built on the two
top-level groups. The \texttt{auth} module registers its login, signup, and
password endpoints on the \texttt{public} group and only its account-deletion
route on \texttt{protected}. Feature modules that belong entirely to one role
create a subgroup and apply a role guard once, as \texttt{exam} does:

\begin{lstlisting}[language=Go,caption={A whole module gated to one role in \texttt{internal/exam/routes.go}.},label={lst:examroutes}]
func RegisterRoutes(public, protected *gin.RouterGroup) {
    exams := protected.Group("/exams")
    exams.Use(middleware.RequireRole("teacher"))
    {
        exams.POST("", CreateExam)
        exams.GET("", GetExams)
        exams.GET("/:id", GetExam)
        exams.PUT("/:id", UpdateExam)
        exams.DELETE("/:id", DeleteExam)
        exams.POST("/:id/assign", AssignExam)
        exams.POST("/:id/close", CloseExam)
        exams.POST("/:id/reopen", ReopenExam)
        exams.GET("/:id/results", GetExamResults)
    }
}
\end{lstlisting}

Here \texttt{/exams} is created on the already-authenticated \texttt{protected}
group and then further restricted so that \emph{every} exam route requires the
teacher role --- two of the three authorization gates applied declaratively
before any handler runs. The \texttt{student} module mirrors this with
\texttt{RequireRole("student")}, and notably reuses the \texttt{exam} package's
attempt handlers, mounting \texttt{exam.StartAttempt},
\texttt{exam.AutosaveAttempt}, and \texttt{exam.SubmitAttempt} under
\texttt{/student/exams/:id/...}: the student module owns those URLs while the
exam module owns the logic, a clean reuse that avoids duplicating attempt code.

The third granularity is a per-route guard, used where endpoints in a module
serve mixed audiences. The \texttt{submission} module opens \texttt{POST /run}
and \texttt{GET /:id} to any authenticated user --- because students run and
read their own submissions and teachers read submissions on exams they own, with
ownership resolved inside the handler --- but pins the grade endpoint to
teachers inline:

\begin{lstlisting}[language=Go,caption={A per-route role guard in \texttt{internal/submission/routes.go}.},label={lst:subroutes}]
subs := protected.Group("/submissions")
{
    subs.POST("/run", RunSolution)
    subs.GET("/:id", GetSubmission)
    subs.PUT("/:id/grade", middleware.RequireRole("teacher"), GradeSubmission)
}
\end{lstlisting}

Passing \texttt{middleware.RequireRole("teacher")} as an argument to a single
route shows that Gin allows role guards to be applied group-wide or
per-endpoint, and APEX uses whichever granularity fits the module.

\subsection{JSON contracts and status codes}

Request and response bodies are plain JSON objects, and the API is disciplined
about status codes because the SPA relies on them. A successful read returns
200 with the resource; a successful create returns the created row. Validation
failures return 400, an absent bearer token or invalid one returns 401, a role
or ownership violation returns 403, and a missing resource returns 404.
Semantically meaningful conflicts get their own codes --- submitting an exam
attempt that is already submitted returns 409 Conflict rather than a generic
error, so the client can react precisely. A subtle but deliberate refinement is
that owner-scoped queries return 404 rather than 403 when a teacher references
another teacher's resource: because the query filters by owner
(\texttt{WHERE id = ? AND teacher\_id = ?}), the row simply ``does not exist''
for the caller, which avoids confirming that the identifier is real. Two further
contract-level choices protect users: login always returns the generic ``invalid
email or password'' regardless of which field was wrong, and
\texttt{POST /auth/forgot-password} always returns 200 --- both to avoid leaking
whether an account exists.

\section{Package Layout and Separation of Concerns}\label{sec:packages}

The backend source lives under \texttt{internal/}, Go's convention for packages
that may not be imported outside the module. Each package owns one slice of the
domain and exposes a small surface --- typically a \texttt{RegisterRoutes}
function and a set of handlers. The organisation is by \emph{feature} rather
than by technical layer: there is no top-level \texttt{controllers/},
\texttt{services/}, and \texttt{models/} trichotomy, but rather a package per
resource that internally spans the transport and handler layers.

\begin{itemize}
  \item \textbf{Cross-cutting infrastructure:} \texttt{config} (environment
        loading and connection-string construction), \texttt{database} (the
        GORM handle, the connection pool, and both migration tracks),
        \texttt{middleware} (CORS, JWT authentication, role gating), and
        \texttt{models} (the fifteen GORM structs that are the single source of
        truth for the schema).
  \item \textbf{Feature modules:} \texttt{auth}, \texttt{class},
        \texttt{student}, \texttt{teacher}, \texttt{exam}, \texttt{problem},
        \texttt{testcase}, \texttt{submission}, \texttt{execute},
        \texttt{folder}, \texttt{announcement}, \texttt{notification}, and
        \texttt{profile}. Each maps to a coherent chunk of the teacher or
        student workflow.
  \item \textbf{Domain services:} \texttt{judge0} (the sandbox client, the
        coding/MCQ/written graders, and the attempt aggregator),
        \texttt{email} (an SMTP wrapper with a stdout fallback), and
        \texttt{reminder} (the one-minute background ticker). These hold logic
        that handlers call into but that is not itself an HTTP endpoint.
\end{itemize}

Alongside \texttt{internal/} sit two command binaries under \texttt{cmd/} ---
\texttt{cmd/migrate} and \texttt{cmd/seed} --- and the root \texttt{main.go} that
wires everything together. The separation of concerns is therefore visible in
the directory tree itself: a reader looking for how exams are assigned opens
\texttt{internal/exam}, and a reader looking for how code is graded opens
\texttt{internal/judge0}, without needing to understand the rest of the system.

\section{Design Patterns and Principles}\label{sec:patterns}

Several established patterns and principles recur through the architecture, and
naming them makes the design intent explicit.

\paragraph{Middleware / chain of responsibility.} The request pipeline of
\cref{fig:layers} is a chain-of-responsibility~\cite{fowler2002patterns}: each
middleware either handles the request (by aborting it) or passes it to the next
link with \texttt{c.Next()}. This is what lets authentication, role checking,
and CORS be composed independently and reordered without touching handlers.

\paragraph{Twelve-factor configuration.} APEX follows the
twelve-factor~\cite{twelvefactor} principle of strict separation of
configuration from code. Every operational setting --- database credentials,
the JWT secret, the Judge0 URL, SMTP details --- is read from the environment
through the \texttt{config} package, never hard-coded, and a \texttt{.env} file
is auto-loaded only as a development convenience. The one setting the process
refuses to start without is \texttt{JWT\_SECRET}: if it is empty, shorter than
32 characters, or the literal placeholder \texttt{dev-secret-change-in-prod},
the process calls \texttt{log.Fatal} and exits. This fail-fast validator makes
``production booted with the development secret'' impossible by construction,
and it is the only fast-fail in the entire boot sequence --- every other missing
setting degrades gracefully.

\paragraph{Stateless services and idempotent operations.} Because the API keeps
no session state, it is horizontally trivial to reason about, and several
operations are deliberately idempotent so that a retry or a refresh cannot cause
harm. Starting an exam attempt, for instance, looks up the attempt for the
(student, exam) pair and returns the existing one if present, backed by the
unique index on \texttt{(user\_id, exam\_id)} so that even a concurrent race
cannot create duplicates. The reminder scheduler records ``already sent'' in
dedicated timestamp columns rather than in memory, which makes it idempotent
across process restarts. These choices reflect the data-intensive-application
discipline~\cite{kleppmann2017ddia} of pushing correctness guarantees down to
the database wherever a purely in-code check would be racy.

\paragraph{Graceful degradation.} Finally, the architecture is built so that
optional dependencies stay optional. Missing SMTP configuration logs email
bodies to stdout; a missing \texttt{GOOGLE\_CLIENT\_ID} makes the Google
sign-in endpoint answer 503 while every other authentication path keeps working.
The system is always runnable with the minimum of one database and one secret,
and each external integration can be switched on independently as an operator is
ready for it.

\medskip

Taken together, these decisions describe an architecture that is intentionally
modest in its ambitions and precise in its execution. The remaining chapters
build on this skeleton: \cref{ch:ch05} details the data model that PostgreSQL
holds, \cref{ch:ch07} follows a submission through the grading engine sketched
here, \cref{ch:ch08} audits the three authorization gates, and \cref{ch:ch10}
shows how the single binary and its companions are packaged and deployed.
